% Pista ana-24j-q3 · Anàlisi
% Procedència: PAU Catalunya 2024 · Convocatòria de juny · Sèrie 1 · Qüestió 3
\documentclass[11pt,a4paper]{article}
\usepackage{pista}
\pistainfo{Catalunya 2024}{Convocatòria de juny}{Sèrie 1}

\begin{document}

\section*{\textcolor{blauPista}{Qüestió 3}}

\noindent En Joan troba entre els papers del seu avi un esbós com el de la figura adjunta, on es descriu un terreny de regadiu que ha deixat en herència al seu pare. La corba de la gràfica és $y = f(x)$, amb $f(x) = -x^{3} + 7x^{2} - 6x + 5$.

\subsection*{\textit{a)} \normalfont A partir de l'expressió de $f(x)$, calculeu les coordenades dels punts $P$, $Q$ i $R$ indicats a la figura. Calculeu també l'equació de la recta $PR$. \hspace*{\fill}\textnormal{[1,25~punts]}}

\pista{Fixa't a la figura que els punts $P$ i $Q$ són sobre la recta horitzontal $y = 5$. Per tant, les seves abscisses són solucions de l'equació
\[
f(x) = 5
\]
Resta $5$ als dos costats: et quedarà una equació sense terme independent, i podràs treure factor comú $x$. Sempre que tens aquesta situació, saps que una solució és $x = 0$, i correspon al punt $(0, 5)$ de la figura. Les altres dues solucions provenen de resoldre una equació de segon grau, i corresponen als punts $P$ i $Q$.

\medskip
A continuació, $R$ és sobre l'eix d'abscisses, just a sota de $Q$: té la mateixa abscissa que $Q$, i ordenada $0$. Finalment, calcula el pendent de la recta $PR$:
\[
m = \frac{y_{R} - y_{P}}{x_{R} - x_{P}}
\]
Després, escriu l'equació de la recta en la forma punt-pendent.}

\subsection*{\textit{b)} \normalfont Calculeu la superfície del terreny. \hspace*{\fill}\textnormal{[1,25~punts]}}

\pista{Fixa't a la figura que la regió ombrejada comença a $P$, on la corba i la recta es tallen, i acaba al segment vertical $QR$. Per tant, integra des de $x_{P}$ i fins a $x_{Q}$. Això vol dir que has de tenir en compte la primera coordenada del punt $P$ i la primera coordenada del punt $Q$, que ja tens de l'apartat \textit{a)}:
\[
\text{Àrea} = \left| \int_{x_{P}}^{x_{Q}} \bigl(f(x) - r(x)\bigr)\,dx \right|
\]
on $r(x)$ és l'equació de la recta $PR$ amb la $y$ aïllada. Fes la resta $f(x) - r(x)$ abans d'integrar. Obtindràs un polinomi de tercer grau, que serà senzill d'integrar.}

\end{document}
